\section*{Scope of this volume}
This volume is the technical reference for
\emph{Generalized Hierarchical Bayesian Segmentation with Irregular Designs,
Multi-Sequence Hierarchies, and Grouped/Latent-Group Designs}. It is based on the editable
manuscript, the implementation archive, and the populated result files supplied with the project.
The statistical statements, computational results, and implementation status are distinguished so
that each conclusion can be read under its stated assumptions.

\begin{statusbox}[Statistical formulation]
\BayesBreak is a generalized hierarchical Bayesian multiple-changepoint method. For regular
exponential-family observation models with proper conjugate priors, segment-specific parameters
can be integrated out to obtain exact segment marginal likelihoods and posterior moments. These
quantities enter sum-product dynamic programming for posterior distributions over ordered
partitions. A separate max-sum recursion with backtracking returns a joint MAP partition. The
method also treats irregular designs, common or group-specific changepoints across multiple
sequences, finite latent groups, and explicitly identified numerical segment integration for
nonconjugate models.
\end{statusbox}

\section*{Status terminology}
\begin{longtable}{@{}p{0.20\textwidth}p{0.70\textwidth}@{}}
\toprule
\textbf{Status} & \textbf{Meaning in this volume}\\ \midrule
\EstablishedTag & A complete proof is given under the stated assumptions.\\
\ExactDerivedTag & The expression follows by direct algebra or exact finite summation.\\
\RealResultTag & A populated value or figure is an archived executed computation.\\
\PartialTag & An executed result addresses only part of a stated question or lacks a required comparison or uncertainty analysis.\\
\ConditionalTag & A conclusion follows only if an additional assumption, such as a uniform numerical error bound, is established.\\
\TargetTag & A proof, implementation correction, or experiment remains to be completed and is not used as an established result.\\
\ExcludedTag & The computation was executed, but its observation model, coordinate axis, reference quantity, or support rule is incompatible with the scientific interpretation under discussion.\\
\bottomrule
\end{longtable}

\begin{exclusionbox}[Archived computations excluded from stated conclusions]
No populated archived numerical value is silently changed. Two computations are retained in the
result history but excluded from the analyses for which they were originally reported:
\resultid{RES-BB-CMP-002}, whose array-CGH comparator used incompatible coordinate axes, and
\resultid{RES-BB-RD-007Q}, whose methylation held-out calculation used a Gaussian predictive
routine for Beta-valued observations and assigned out-of-range coordinates to the last segment.
A corrected calculation must receive a new result identifier.
\end{exclusionbox}

\section*{Reading routes}
\begin{itemize}
\item \textbf{Statistical formulation and exact inference:} Chapters~\ref{ch:orientation},
\ref{ch:formal}, \ref{ch:block}, \ref{ch:sumproduct}, and \ref{ch:map}.
\item \textbf{Irregular designs and hierarchical extensions:} Chapters~\ref{ch:priors},
\ref{ch:shared}, \ref{ch:templates}, and \ref{ch:nonconj}.
\item \textbf{Prediction, implementation, and empirical studies:} Chapters~\ref{ch:prediction},
\ref{ch:software}, \ref{ch:synthetic}, \ref{ch:realdata}, and
Appendices~\ref{app:implementation}--\ref{app:traceability}.
\item \textbf{Literature and limitations:} Chapters~\ref{ch:lineage}, \ref{ch:failures}, and
Appendix~\ref{app:annotated-lit}.
\end{itemize}

\section*{Source precedence}
The editable manuscript and its archived numerical assets are the factual baseline. The author's
explicit decisions fix the title and main narrative. Mathematical corrections are stated in place;
implementation defects and unresolved experiments are reported without rewriting the archived
outputs. External results are attributed to their sources and are not presented as BayesBreak
results.
